\documentclass[conference]{IEEEtran}
\IEEEoverridecommandlockouts
\usepackage{cite}
\usepackage{amsmath,amssymb,amsfonts}
\usepackage{algorithmic}
\usepackage{graphicx}
\usepackage{textcomp}
\usepackage{xcolor}
% URL package to reference urls in bib without getting "Missing $ inserted" error and "LaTeX: Label(s) may have changed. Rerun to get cross-references right.LaTeX" warning.
\usepackage{url}
\usepackage{tabularx}
\usepackage{array}
\newcolumntype{C}[1]{>{\centering\arraybackslash\hyphenpenalty=10000}p{#1}}
\newcolumntype{Y}{>{\centering\arraybackslash\hyphenpenalty=10000}X}
\begin{document}

\title{A Comparative Analysis of File System Type and Digital Forensic Data Recoverability
\\
}

\author{\IEEEauthorblockN{Ahmed Ali, Hanan Kakooza, Evan Vink, Royce Coviello}
\IEEEauthorblockA{Team 1: DDR4. Digital Forensics and Cybersecurity\\
 British Columbia Institute of Technology (BCIT)\\
 Vancouver, BC, Canada\\
aali148@my.bcit.ca, hkakooza1@my.bcit.ca, evink2@my.bcit.ca, rcoviello@my.bcit.ca}
}

\maketitle
\begin{abstract}

Digital forensic investigations often require recovering deleted or reformatted data from storage media. Existing research has examined file-carving techniques, recovery tools, and individual file systems. However, comparatively less research has directly evaluated how different file-system architectures affect data recoverability under identical experimental conditions. This study proposes a controlled, reproducible comparison of forensic file recovery across NTFS, exFAT, ext4, Btrfs, and bcachefs using a standardized NixOS virtual-machine environment.

This study will experimentally evaluate whether file-system architecture affects forensic data recoverability under controlled conditions. A reproducible NixOS virtual-machine environment will be used to test NTFS, exFAT, ext4, Btrfs, and bcachefs using fresh virtual disks for each experimental run. A standardized dataset containing known files will be written to each file system and cryptographic hashes will be recorded before a controlled data-loss procedure is performed and file recovery will then be attempted using PhotoRec.

This study aims to determine whether file-system type produces measurable differences in forensic data recoverability when hardware-related factors are intentionally controlled. The resulting virtual test environment is intended to provide a reproducible baseline for comparing file-system behaviour.


\end{abstract}

\begin{IEEEkeywords}
component, formatting, style, styling, insert
\end{IEEEkeywords}

\section{Introduction}

\subsection{Background and Context}

Even after users delete files, format storage media, or lose access to data, digital storage devices often retain substantial residual information that remains valuable for forensic investigations \cite{liambas2024performance}. Sibe and Idigbo \cite{sibe2025pii} highlight real-world consequences of persistent residual data by examining refurbished hard drives that had undergone formatting and sanitization, successfully retrieving 28,691 files which included sensitive personally identifiable information.

File carving is used when file-system metadata is corrupted, incomplete, or missing, as it reconstructs file content directly from unallocated space using recognized file signatures and internal structural patterns. Tam and Abdullah \cite{tam2024androidjpeg} highlight PhotoRec as an effective carving utility that functions independently of active file-system structures and supports numerous file formats. In an experimental evaluation within the CAINE forensic environment, Carranza and DeCusatis \cite{carranza2021daubert} demonstrated that PhotoRec yields consistent, repeatable recovery performance, reinforcing its suitability as a standardized open-source carving tool. Further empirical studies confirm that PhotoRec can effectively carve target storage partitions without requiring a forensic image. For instance, Pratama \cite{pratama2021photorec} applied PhotoRec directly to a FAT32 USB flash drive and successfully recovered all 2,781 test files, while Dasmen et al.\ \cite{dasmen2024photorec} demonstrated the successful retrieval of diverse file types from external flash media. While these findings do not guarantee universal file recovery across all scenarios, they confirm PhotoRec's utility as a reliable signature-based recovery tool.

Nevertheless, recovery success is heavily influenced by the specific file system and operational environment under investigation. Researching NTFS, Venkatesh et al.\ \cite{venkatesh2024ntfs} achieved an approximate 95\% recovery rate by parsing Master File Table (MFT) records to restore deleted files. In Linux environments, Kim et al.\ \cite{kim2021ext4xfs} demonstrated that alterations to ext4 journal structures and XFS metadata directly affect tool performance and hash integrity during file reconstruction. Btrfs introduces unique challenges, as Pandey et al.\ \cite{pandey2026btrfs} note its copy-on-write (CoW) design, subvolumes, and historical B-tree roots preserve recoverable data states but require distinct reconstruction methodologies compared to traditional file systems. Collectively, these insights show that file-system architecture is a key aspect of forensic recoverability and should be evaluated as an explicit experimental variable.

Rather than directly comparing underlying file systems, many other studies focus on specialized embedded platforms. For example, Raymer et al.\ \cite{raymer2023quest2} established a standardized forensic acquisition and analysis protocol for the Meta Quest 2, whereas Marques et al.\ \cite{marques2026bmw} investigated QNX6FS on BMW infotainment systems, utilizing PhotoRec to carve unallocated space. Although these investigations emphasize the importance of rigorous, reproducible methodologies, they also illustrate how current literature remains largely tied to specific hardware devices, specialized operating systems, or isolated file-system implementations.

\subsection{Problem Statement}

File recovery and file carving are established areas of digital forensics, but recovery outcomes may depend not only on the recovery tool but also on the underlying file system. File systems use different structures to allocate data, record metadata, track deletion, and manage free space. In NTFS, for example, deleted file entries may remain represented in the MFT while their data persists until the associated storage space is reused \cite{venkatesh2024ntfs}. Ext4 can retain deleted-file metadata within its journal \cite{kim2021ext4xfs}, whereas Btrfs uses CoW B-trees and historical filesystem roots rather than traditional fixed metadata structures \cite{pandey2026btrfs}. These architectural differences create a reasonable basis for expecting file-system type to affect what remains recoverable after deletion.

Existing studies provide useful evidence about individual recovery environments, but they generally do not isolate filesystem type under the same recovery environment. Chhajed \cite{chhajed2024persistence} is particularly relevant because the study compares deleted-file persistence across several file systems on both HDD and SSD storage, however, it also varies by operating systems, disk types, and ordinary user activities, making it difficult to attribute observed differences to filesystem architecture alone. Kumar and Haribabu \cite{kumar2022docker} hold a Linux partition constant and compare three carving tools against deleted Docker-container data, while Venkatesh et al.\ \cite{venkatesh2024ntfs}, Kim et al.\ \cite{kim2021ext4xfs}, and Pandey et al.\ \cite{pandey2026btrfs} focus on filesystem-specific recovery techniques. The literature therefore demonstrates that both recovery method and filesystem structure matter, but leaves room for a controlled experiment in which the recovery tool and virtual environment are held constant while only the filesystem is changed.

This study addresses this problem by creating a controlled virtual test environment in which NTFS, exFAT, ext4, Btrfs, and bcachefs undergo the same data-generation, deletion, and recovery procedures. By using fresh virtual disks for each trial, the study seeks to establish a reproducible software-level baseline of file-system recoverability without introducing hardware-specific behaviour such as physical media wear or storage-controller effects.

\subsection{Literature Overview}

Existing work relevant to this study can be grouped into four broad themes. First, researchers have evaluated forensic tools and their performance. Carranza and DeCusatis \cite{carranza2021daubert} experimentally validated open-source tools including PhotoRec within CAINE, while Liambas and Manios \cite{liambas2024performance} compared PhotoRec and other forensic utilities across multiple operating systems and repeated each tool test, ten times to obtain the average execution times. Pratama \cite{pratama2021photorec} and Dasmen et al.\ \cite{dasmen2024photorec} both used PhotoRec for direct recovery from removable storage. Kumar and Haribabu \cite{kumar2022docker} compared Scalpel, Foremost, and PhotoRec when carving data from deleted Docker containers, providing a useful precedent for controlled recovery experiments conducted in virtualized or software-defined environments.

Second, several studies investigate file recovery within a particular file system or storage architecture. Venkatesh et al.\ \cite{venkatesh2024ntfs} developed a Python/PyTSK3 approach for recovering deleted NTFS files through MFT analysis. Kim et al.\ \cite{kim2021ext4xfs} developed a TSK-based framework that uses ext4 journal metadata for deleted-file recovery and XFS metadata for extraction, showing that forensic tool support must adapt when filesystem metadata structures change. Choi and Lee \cite{choi2022raid} similarly combine NTFS metadata analysis with file carving to recover data from damaged RAID systems. Most directly relevant to the present filesystem set, Pandey et al.\ \cite{pandey2026btrfs} show that Btrfs recovery is shaped by CoW B-trees, historical roots, subvolumes, and that a filesystem-aware method can reconstruct the file content, paths, and metadata when the required structures survive.

Third, file-carving studies have emphasized content-based recovery and integrity verification. Tam and Abdullah \cite{tam2024androidjpeg} describe file carving as recovery from raw or unallocated space without depending on active file-system metadata, and their Android JPEG recovery workflow uses hash values to verify recovered-file integrity. Kumar and Haribabu \cite{kumar2022docker} found that their carving tools did not preserve original filenames or timestamps and that recovery success varied by file type in their Docker dataset. Pandey et al.\ \cite{pandey2026btrfs} also distinguish signature-based carving from filesystem-aware reconstruction, noting that carving generally provides weaker path, metadata, fragmentation, and filesystem-consistency guarantees. These findings support the present study's decision to identify recovered artifacts by cryptographic hash rather than filename and to treat exact content integrity as a primary outcome.

Finally, research on real storage devices demonstrates why environmental control is necessary. Sibe and Idigbo \cite{sibe2025pii} recovered substantial residual data from refurbished hard drives that had reportedly been formatted and sanitized. Chhajed \cite{chhajed2024persistence} found that deleted-file persistence can be affected by the filesystem, storage medium, operating system, and subsequent user activity. These studies provide strong real-world motivation, but they also illustrate the confounding factors introduced by physical HDDs and SSDs. The present study intentionally removes those hardware variables so that filesystem-dependent differences in PhotoRec recovery can be examined against a common virtual baseline.

\subsection{Research Gap}

Taken together, the initial literature establishes that deleted-data recovery, file carving, forensic-tool performance, and file-system-specific recovery are active research areas. However, the studies reviewed so far generally examine a single file system, a particular device or storage medium, or differences between recovery tools and operating environments. None of the articles in this first set directly compare several modern file systems under the same virtual hardware, test corpus, deletion procedure, and PhotoRec recovery process.

The proposed study therefore shifts the file system itself to the primary independent variable. By testing NTFS, exFAT, ext4, Btrfs, and bcachefs in fresh, standardized virtual disks, the experiment is designed to reduce hardware-related confounders and isolate differences in logical recoverability.

The research gap addressed in this proposal is therefore the limited availability of controlled and repeatable experiments that compare forensic file recovery across multiple file-system architectures while keeping the recovery tool and experimental environment constant.

\subsection{Research Objectives}

The objective of this research is to deploy an automated, repeatable virtual testing environment to compare forensic data-recovery outcomes across different file systems while minimizing variation from hardware and experimental setup. This study aims to explore how the underlying filesystem architecture (NTFS, ext4, exFAT, Btrfs, and bcachefs) affects the subsequent file carving through PhotoRec.

\section{Literature Review}

The literature relevant to this study spans two related but distinct approaches to deleted-data recovery. The first is \emph{content-based file carving}, in which tools such as PhotoRec search storage for recognizable file signatures and structures without depending on an intact filesystem directory. The second is \emph{filesystem-aware recovery}, in which investigators interpret filesystem-specific metadata such as the NTFS Master File Table (MFT), Ext4 journals, FAT allocation structures, or Btrfs B-trees. Both approaches show that deleted data can remain recoverable, but they measure different aspects of recoverability. For this reason, the present study keeps PhotoRec constant and treats filesystem type as the independent variable rather than attempting to compare fundamentally different recovery algorithms.

We organize the reviewed work into three themes: (1) file-carving tools and controlled recovery experiments, (2) filesystem-specific recovery and metadata structures, and (3) storage-medium effects, virtualization, and reproducible forensic testing. Table~\ref{tab:lit_matrix} summarizes representative studies and the specific limitation each leaves for the present experiment.

\subsection{Literature Review Matrix}

\begin{table*}[t]
\caption{Literature Review Matrix}
\label{tab:lit_matrix}
\centering
\scriptsize
\setlength{\tabcolsep}{2.7pt}
\renewcommand{\arraystretch}{1.14}
\begin{tabularx}{\textwidth}{|p{2.15cm}|C{0.65cm}|>{\raggedright\arraybackslash}X|>{\raggedright\arraybackslash}X|>{\raggedright\arraybackslash}X|>{\raggedright\arraybackslash}X|}
\hline
\textbf{Study} & \textbf{Year} & \textbf{Method / Scope} & \textbf{Main Finding} & \textbf{Relevance to This Study} & \textbf{Remaining Limitation / Gap} \\
\hline
Carranza and DeCusatis \cite{carranza2021daubert} & 2021 & Repeated validation of open-source forensic tools in the CAINE environment, including PhotoRec. & PhotoRec produced consistent recovery behaviour across repeated scans. & Supports using a documented open-source carver in a repeatable experimental workflow. & Validates tools rather than comparing filesystem architectures. \\
\hline
Pratama \cite{pratama2021photorec} & 2021 & Direct PhotoRec recovery of 2,781 erased files from a FAT32 USB flash drive. & All 2,781 test files were reported recovered in the tested scenario. & Provides precedent for direct source-partition scanning and quantitative recovery counts. & One filesystem and one physical medium, does not isolate filesystem effects. \\
\hline
Kumar and Haribabu \cite{kumar2022docker} & 2022 & Deleted Docker-container data carved with Scalpel, Foremost, and PhotoRec using multiple file types. & Recovery varied by tool and file type, PhotoRec recovered 5/10 files and original names/timestamps were not preserved. & Closely parallels a controlled software-defined recovery environment and supports hash-based identification rather than filenames. & Compares carving tools, not filesystems, small test corpus. \\
\hline
Tam and Abdullah \cite{tam2024androidjpeg} & 2024 & JPEG carving and geotag analysis on Android images with integrity verification. & Recovered JPEGs were checked using hash comparison in addition to metadata analysis. & Supports cryptographic hash verification as a measure of exact recovery. & Focused on JPEG/mobile forensics rather than cross-filesystem comparison. \\
\hline
Liambas and Manios \cite{liambas2024performance} & 2024 & Performance comparison of forensic tools across operating systems with repeated trials. & Repeated testing was used to calculate average execution performance. & Supports repeated runs and recording processing time rather than relying on a single trial. & Varies operating environments and tools rather than isolating filesystem type. \\
\hline
Venkatesh et al. \cite{venkatesh2024ntfs} & 2024 & Deleted-file recovery from NTFS using Python, PyTSK3, and MFT metadata. & Demonstrated high recovery in the test images and showed the forensic value of surviving MFT records. & Shows directly that filesystem-specific metadata can influence recoverability after deletion. & NTFS-specific, metadata-aware method, not a generic carving comparison. \\
\hline
Kim et al. \cite{kim2021ext4xfs} & 2021 & TSK-based Ext4 journal recovery and XFS metadata extraction, compared with existing tools. & Metadata/version changes affected recovery support and integrity, the proposed filesystem-aware framework improved recovery accuracy. & Demonstrates that filesystem design and version-specific metadata can change recovery outcomes. & Limited to Ext4/XFS and uses filesystem-aware parsing rather than one common carver. \\
\hline
Kuts et al. \cite{kuts2023fat32} & 2023 & Experimental deleted-file recovery in FAT32 using its cluster and directory structures. & Recovery success varied by file type and surviving cluster information, signature recovery could lose metadata. & Shows that deletion semantics and file layout can influence what a recovery tool reconstructs. & FAT32-specific and does not test the modern filesystem set used here. \\
\hline
Pandey et al. \cite{pandey2026btrfs} & 2026 & Deterministic deleted-file recovery by parsing Btrfs CoW trees, historical roots, inodes, and extents. & Btrfs historical structures can support content, path, and metadata reconstruction when required structures survive. & Provides strong evidence that CoW filesystem architecture creates recovery behaviour unlike traditional filesystems. & Btrfs-specific specialized algorithm, its results are not directly equivalent to PhotoRec carving. \\
\hline
Aigbogun and Zhou \cite{aigbogun2025f2fs} & 2025 & Forensic study of F2FS and emerging file systems. & Highlights the need for forensic methods to evolve with newer filesystem designs and metadata structures. & Supports including modern filesystems rather than limiting comparisons to legacy FAT/NTFS families. & Does not provide a controlled PhotoRec comparison across the filesystems in this experiment. \\
\hline
Chhajed \cite{chhajed2024persistence} & 2024 & Compared persistence of deleted files across different filesystems and HDD/SSD media. & Persistence was affected by filesystem, storage medium, operating environment, and subsequent activity. & One of the closest comparative studies and motivates treating filesystem as an explicit variable. & Hardware, operating-system, and activity effects make filesystem-only attribution difficult. \\
\hline
Sibe and Idigbo \cite{sibe2025pii} & 2025 & Recovery of residual data from three refurbished hard drives reportedly formatted/sanitized. & One drive yielded 28,691 recoverable files including sensitive PII. & Demonstrates the practical forensic and privacy importance of residual deleted data. & Real devices introduce uncontrolled history, hardware state, and sanitization differences. \\
\hline
Iqbal and Soewito \cite{iqbal2020ssdtrim} & 2020 & Static forensic recovery from SSD storage with TRIM and Deep Freeze across repeated scenarios. & Recovery became less successful across later scenarios, demonstrating the effect of SSD behaviour on residual data. & Supports separating logical filesystem effects from physical-storage mechanisms in the first phase of the project. & Hardware/firmware behaviour is central, so filesystem effects cannot be isolated. \\
\hline
Chen et al. \cite{chen2022ransomware} & 2022 & Ransomware-data recovery using filesystem forensics together with flash translation layer extraction. & Demonstrates that evidence can exist below the logical filesystem layer. & Motivates treating hardware/FTL recovery as a future extension rather than mixing it into the virtual baseline. & Focused on ransomware and flash-layer recovery rather than standardized filesystem comparison. \\
\hline
Wickramasekara et al. \cite{wickramasekara2026autodfbench} & 2026 & Benchmarking framework for repeatable digital-forensic tool testing and generated-code evaluation. & Emphasizes standardized, reproducible test cases and measurable outputs when evaluating forensic tools. & Supports automation, fixed test data, repeated trials, and structured result collection. & General benchmarking framework rather than a study of deleted-file recovery across filesystems. \\
\hline
\end{tabularx}
\end{table*}

\subsection{Theme 1: File Carving and Controlled Recovery Experiments}

A substantial portion of the literature evaluates whether deleted content can be reconstructed using general-purpose carving tools. PhotoRec is especially relevant because it searches the underlying storage for recognizable file structures rather than depending on active filesystem directory entries. Carranza and DeCusatis \cite{carranza2021daubert} validated PhotoRec as part of an open-source CAINE forensic workflow and reported consistent behaviour over repeated scans. Pratama \cite{pratama2021photorec} likewise used PhotoRec directly against a FAT32 USB flash drive and reported recovery of all 2,781 test files. Dasmen et al. \cite{dasmen2024photorec} further demonstrate PhotoRec's use for recovering multiple file types from external storage. These studies establish PhotoRec as an appropriate controlled recovery tool, but their high recovery rates should not be generalized to every deletion scenario, filesystem, or storage medium.

Other research shows why a standardized dataset and integrity test are necessary. Kumar and Haribabu \cite{kumar2022docker} populated Docker containers with several file types, deleted the containers, and carved the underlying partition using Scalpel, Foremost, and PhotoRec. Their recovery rates differed across tools, with PhotoRec recovering 5 of 10 files, while original filenames and timestamps were not preserved. This result is important for the present experiment because a recovered filename cannot be assumed to identify the original artifact. Instead, the proposed testbed compares recovered content against pre-deletion SHA-256 hashes. Tam and Abdullah \cite{tam2024androidjpeg} similarly use hashing to verify recovered JPEG integrity, providing precedent for treating exact cryptographic identity as a stronger outcome than simple file detection.

The Docker study is also methodologically useful because it performs carving directly against the relevant partition rather than requiring an intermediate forensic image \cite{kumar2022docker}. This resembles the proposed NixOS testbed, where the virtual test disk is unmounted and then scanned directly by PhotoRec. The purpose of the current experiment is therefore not to evaluate forensic acquisition procedures, but to hold the recovery procedure constant while examining whether the logical filesystem environment changes what PhotoRec can reconstruct. Liambas and Manios \cite{liambas2024performance} further support repeated benchmarking rather than one-off measurements, while AutoDFBench \cite{wickramasekara2026autodfbench} provides broader support for standardized test cases and reproducible evaluation of forensic tools.

\subsection{Theme 2: Filesystem Architecture and Deleted-Data Recovery}

Filesystem-aware research demonstrates that deletion is not a uniform operation across storage formats. NTFS, Ext4, FAT32, XFS, Btrfs, and newer Linux filesystems maintain different allocation and metadata structures, so the forensic traces that survive deletion can differ substantially. Venkatesh et al. \cite{venkatesh2024ntfs} recover deleted NTFS files by interpreting MFT records, illustrating that metadata may continue to describe a deleted file until relevant records or data regions are reused. Zhang et al. \cite{zhang2020windows} similarly compare recovery logic for FAT32 and NTFS and attribute differences in recovery to their distinct storage and metadata mechanisms.

Linux filesystems introduce different sources of recoverable state. Kim et al. \cite{kim2021ext4xfs} show that Ext4 journal structures and XFS metadata versions affect extraction and deleted-file recovery, and that existing tools may produce recovered files whose hashes do not match the originals. Their results reinforce the need to distinguish a tool merely outputting a file from an exact, verified recovery. They also show that forensic support can lag behind filesystem changes, because a tool designed around one metadata version may not correctly interpret a newer structure. Aigbogun and Zhou \cite{aigbogun2025f2fs} extend this concern to F2FS and other emerging filesystems, showing why contemporary filesystem designs remain an active forensic research problem.

Btrfs is especially important to the present study because its copy-on-write architecture differs substantially from traditional allocation-table or journal-based designs. Pandey et al. \cite{pandey2026btrfs} reconstruct deleted files by traversing Btrfs B-trees and historical roots, recovering not only content but also paths and inode metadata when the necessary structures remain resolvable. Their work also identifies an important limitation of signature carving: a carver may recover recognizable file content without reconstructing filesystem provenance, directory structure, timestamps, or other metadata. Therefore, this project does not claim to measure the \emph{maximum possible} recovery from Btrfs, NTFS, Ext4, or any other filesystem. It measures recovery through one standardized carving method, PhotoRec, so that the same recovery process is applied to every filesystem.

The reviewed literature is comparatively strong for NTFS, FAT32, Ext4, XFS, and Btrfs, but the current corpus contains much less direct experimental work on exFAT and especially bcachefs. This does not establish that such research does not exist, rather, it identifies an imbalance within the literature collected for this project and strengthens the value of including newer filesystems in a common experimental baseline.

\subsection{Theme 3: Storage Media, Virtualization, and Experimental Control}

Physical storage behaviour is a major confounding factor in deleted-data recovery. Sibe and Idigbo \cite{sibe2025pii} recovered 28,691 files containing sensitive information from one of three refurbished drives that had reportedly been formatted and sanitized, illustrating both the practical importance and the unpredictability of residual data on real devices. Chhajed \cite{chhajed2024persistence} is particularly close to the present research question because it compares deleted-file persistence across multiple filesystems and disk types. However, its results also reflect storage medium, operating-system behaviour, and subsequent user activity, making it difficult to attribute differences exclusively to filesystem architecture.

SSD-specific studies make this problem more explicit. Iqbal and Soewito \cite{iqbal2020ssdtrim} examine recovery in an SSD environment involving TRIM and show that recoverability declines across their scenarios. Chen et al. \cite{chen2022ransomware} investigate recovery using both filesystem forensics and flash translation layer data extraction, demonstrating that relevant remnants may exist below the logical filesystem. These studies are important for real-world forensics, but they also show why introducing physical flash hardware in the present experiment would make the independent variable harder to isolate.

The NixOS virtual-machine design is therefore intended as a controlled first-stage model rather than a complete reproduction of physical storage. Each trial uses a fresh virtual disk, the same test corpus, the same deletion and overwrite procedure, the same PhotoRec version, and the same recovery destination. Virtualization cannot model controller firmware, TRIM implementation, garbage collection, wear levelling, or physical media degradation, but excluding those mechanisms is deliberate: it allows the experiment to establish a software-level baseline before hardware effects are introduced in future work. The controlled software environment also follows the broader emphasis on repeatability seen in Docker-based carving experiments \cite{kumar2022docker} and forensic benchmarking work \cite{wickramasekara2026autodfbench}.

\subsection{Synthesis and Research Gap}

The literature establishes three points relevant to the proposed experiment. First, signature-based carving is a viable recovery method, but its success varies with file type, recovery tool, and test conditions \cite{pratama2021photorec,kumar2022docker}. Second, filesystem-specific studies show that metadata organization, journaling, allocation, and copy-on-write behaviour directly affect which traces remain after deletion \cite{venkatesh2024ntfs,kim2021ext4xfs,pandey2026btrfs}. Third, physical storage mechanisms can independently change data persistence, particularly on SSD and flash-based media \cite{iqbal2020ssdtrim,chen2022ransomware,chhajed2024persistence}.

What remains limited in the reviewed work is a controlled comparison in which several modern filesystems are subjected to the same test corpus, deletion/overwrite procedure, virtual hardware, software versions, and PhotoRec recovery process. Existing comparative studies commonly vary more than one of these factors, while filesystem-specific studies frequently use specialized recovery methods that are not directly comparable with one another. The present study addresses that methodological gap by fixing the recovery tool and experimental environment while varying NTFS, exFAT, Ext4, Btrfs, and bcachefs. The resulting measurements should therefore be interpreted as comparative \emph{PhotoRec recoverability under standardized virtual conditions}, not as a universal ranking of the absolute forensic recoverability of each filesystem.

\section{Methodology}
% Explain the comprehensive approach you used to conduct this research here.
This research will use a quantitative, experimental, comparative methodology. It is quantitative because forensic recoverability will be measured using numerical outcomes, including the number and proportion of files recovered, SHA-256 hash matches, and image-content similarity. It is experimental because the researchers will deliberately vary the file system while controlling the virtual storage environment, dataset, recovery software, and data-loss procedure. It is comparative because the same experimental process will be repeated across multiple file systems and the resulting recovery measurements will be compared.

\subsection{Research Method}

The primary independent variable will be file-system type, with five experimental levels: NTFS, exFAT, ext4, Btrfs, and bcachefs. If multiple data-loss scenarios are tested, the data-loss condition will act as a secondary independent variable.
We are specifically using NixOS which will allow us to control specific packages.
The primary dependent variable will be forensic recovery success. This will be measured through exact file recovery, determined by comparison with SHA-256 hashes recorded before deletion. For recoverable image files that fail exact hash verification, secondary measurements will evaluate image dimensions and pixel-level similarity to determine the extent to which visual content was recovered.
Controlled variables will include the virtual-machine configuration, virtual-disk size, source dataset, filesystem preparation process, recovery software, recovery destination, and experimental procedure. A fresh virtual disk will be provided for each test run to prevent residual data from previous experiments from contaminating later trials.
The virtual environment is particularly appropriate for this research because it isolates filesystem behaviour from physical-device characteristics. Consequently, the results should be interpreted as a controlled theoretical baseline rather than a direct simulation of real-world storage behaviour.

\subsection {Data Sampling}

Experiments will be conducted within a standardized NixOS virtual-machine environment. Each VM boot will provide a fresh virtual disk for the filesystem under test and a separate virtual disk for recovery output. This design prevents recovered files from being written back to the source disk and reduces contamination between experimental runs. The experiment involves scanning the unmounted virtual test disk directly with PhotoRec rather than traditionally creating a forensic image. Direct source-disk or partition scanning has precedent in prior PhotoRec and Docker-carving experiments Pratama \cite{pratama2021photorec} and Kumar and Haribabu \cite{kumar2022docker}. However, for this experiment we will be using a virtual environment.

For each trial, the test disk will be formatted using one of the selected file systems: NTFS, exFAT, ext4, Btrfs, or bcachefs. An identical standardized dataset will then be written to the file system. Before deletion, SHA-256 hashes will be generated and stored for later integrity verification. The current recovery environment already follows this model by creating known files and recording their hashes before the data-loss operation.

The data loss scenario will test recovery after a logical delete via the GNU coreutils `rm` command. Next, a different fixed dataset will be written to the recovery target disk to simulate data being overwritten by normal computer use. Afterwards, the disk will be unmounted and PhotoRec will be run against the recovery target disk.

Image files will be subject to an additional evaluation for partially recovered data. The team's current Python implementation generates deterministic 1920×1080-pixel solid-colour PNG images with preset metadata and is capable of comparing images at the individual-pixel level. The analysis code compares image content hashes and when they differ, records the coordinates of pixels with different RGB values.

For each filesystem, the experimental procedure will be as follows:
\begin{enumerate}
    \item A fresh virtual test disk will be created.
    \item The virtual disk will be formatted using the selected filesystem.
    \item The standardized test corpus will be written to the filesystem.
    \item SHA-256 hashes of all original test files will be recorded before the data-loss operation.
    \item A controlled data-loss condition will be applied.
    \item The test filesystem will be unmounted before recovery.
    \item PhotoRec will scan the test disk and write recovered files to the separate recovery-output disk.
    \item Recovered files will be compared with the original SHA-256 log.
    \item Recoverable image files that do not match any hashes will undergo secondary image-content analysis.
    \item The resulting measurements will be written to structured experiment logs for later statistical analysis.
\end{enumerate}

Steps 1-10 will be repeated 10 times for each file system. The average of the test results will be used to determine the file-recoverability of each file system.

\subsection {Data Analysis Techniques}

Recovery results will be analyzed primarily using exact file integrity and recovery-rate measurements.

PhotoRec mainly recovers files by patterns in the file contents rather than file system records. Therefore, the differences in recovery are not solely reliant on the file system architecture. However, with our reproducible VM environment testing each filesystem with the same recovery files, we can control other factors and compare the filesystems in the exact same conditions.

First, each recovered file will be hashed and compared with the SHA-256 hashes recorded in the original log. A matching hash indicates the file was recovered byte-for-byte. For each filesystem and experimental condition, the following values will be calculated:

\textbf{File recovery rate:}
% Proper way to input an equation
\begin{equation}
% shrink the equation to 90% of the column width. --> (0.90\columnwidth)
\resizebox{0.90\columnwidth}{!}{$\displaystyle
\text{File Recovery Rate} = \frac{\text{Number of target files recovered}}{\text{Total number of target files}} \times 100
$}
\end{equation}

\textbf{Exact recovery rate:}
\begin{equation}
\resizebox{0.95\columnwidth}{!}{$\displaystyle
\text{Exact Recovery Rate} = \frac{\text{Recovered files with matching SHA-256 hashes}}{\text{Total number of recovered target files}} \times 100
$}
\end{equation}

\subsection {Data Measurement Models}

We will be using a continuous measurement model to test data recoverability. A matching SHA-256 hash will indicate that the recovered artifact is identical to its baseline version. The existing Nix recovery checker already follows this concept by matching recovered files against the previously generated manifest based on SHA-256 rather than relying on filenames, which may be lost during file carving.
For image artifacts that are recovered but fail integrity verification, additional content-level analysis will be performed. The Python image-analysis program is designed to compare image dimensions, image-content hashes, metadata, and individual pixel values. A pixel-integrity measurement can therefore be expressed as:

\begin{equation}
\resizebox{0.80\columnwidth}{!}{$\displaystyle
\text{Pixel Integrity} = \frac{\text{Matching pixels}}{\text{Total pixels}} \times 100
$}
\end{equation}

This provides a more granular measurement than classifying every non-identical file simply as "unrecovered." Choi and Lee \cite{choi2022raid} note that recovered artifacts can themselves be corrupted when portions of a file have been overwritten, meaning that successful identification of an artifact does not necessarily imply complete content recovery. For example, an image that is 99.9\% visually intact but contains corruption in a small region should be distinguishable from an image that could not be recovered at all.
Recovery speed will also be evaluated. Acquisition time, PhotoRec carving time, and total recovery-processing time should be stored separately. Results for each filesystem will then be summarized using descriptive statistics such as mean, median, standard deviation, minimum, and maximum recovery rate and processing time.
The results can be presented using bar charts for mean recovery rates, box plots for the distribution of recovery times across repeated trials, and a comparison table and heatmap showing exact and partial recovery performance for each filesystem.

\subsection{Timeline \& Milestones}
\begin{table}[h]
\centering
\small
\renewcommand{\arraystretch}{1.3}
\begin{tabularx}{\columnwidth}{|C{1.9cm}|C{2.1cm}|Y|}
\hline
\textbf{Task} & \textbf{Estimated Duration} & \textbf{Milestones} \\
\hline
\textbf{Proposal Submission} & Sept 20 -- Oct 6, 2026 & Identified research gaps, identified research methodology, and problem statement \\
\hline
\textbf{Testbed Setup} & October 7 -- October 20 2026 & Automated VM spin up, hashes, and results \\
\hline
\textbf{Creating Python Scripts} & October 7 -- October 10 2026 & Generate test photos, compare the original generated photo with the recovered photo using automation, and dump results into a log file for data analysis at a later date, as well as a Python log file digester script \\
\hline
\textbf{Data Collection} & October 21 -- October 31 2026 & Finalize testbed and Python scripts for accurate data collection and run experiments to generate data in the log file \\
\hline
\textbf{Data Analysis} & November 1 -- November 5 2026 & Parse and organize experimental results, calculate recovery rates from log files, identify trends and anomalies \\
\hline
\textbf{Draft Writing} & November 6 -- November 30 2026 & Complete Draft, methodology, results, conclusion completed, Included visual tables and graphs for the results \\
\hline
\textbf{Final Report \& Presentation} & December 2026 & Complete final report, Complete final presentation and show our findings \\
\hline
\end{tabularx}
\end{table}



\section{Evaluation}
\subsection{Research Questions/Hypothesis}

The main question behind our research is how different filesystem architectures (NTFS, Ext4, exFAT, Btrfs, and bcachefs) affect the outcome of file carving through PhotoRec.
Is there a significant relationship between different file systems and the ability to recover deleted data using PhotoRec?

\section{Conclusion and Future Work}
What are the existing gaps and unanswered research questions? what is the next step?



\subsection{Future Work}
The present experiment is intentionally designed as a controlled software-level baseline. By using fresh virtual disks in a reproducible NixOS environment, the study reduces hardware-specific effects and allows differences in PhotoRec recovery to be compared while the test corpus, deletion procedure, overwrite workload, and recovery software remain constant. In this sense, the virtual environment represents a theoretical best-case or control condition for logical file recovery rather than a complete representation of field conditions.

A natural extension of this work is to repeat the same experimental procedure on physical USB flash drives and compare the results directly with the virtual baseline. Prior studies show that recovery from physical storage can be affected by the storage medium and subsequent device activity \cite{chhajed2024persistence}, while flash-based storage introduces additional controller-level behavior such as wear levelling, garbage collection, and flash-translation-layer management that may alter the persistence of deleted data \cite{iqbal2020ssdtrim,chen2022ransomware}. Physical USB recovery with PhotoRec has also been demonstrated in prior work \cite{pratama2021photorec}. Future testing would therefore use the same filesystems, test corpus, deletion/overwrite procedure, PhotoRec configuration, and SHA-256 verification used in the virtual experiment, while introducing physical USB hardware as the primary new variable. Comparing several USB devices or models against the VM control would help determine how much real storage hardware changes recovery success, integrity, and processing time.


\section*{Acknowledgment}

\section*{References}
Provide your citations as long as you continue with your research. For example, Einsteign et.al. \cite{einstein1905} demonstrated that energy converts...

% Add all references (without explicitly calling them)
\nocite{*}
\bibliographystyle{IEEEtran}  % IEEE style
\bibliography{refs}

\end{document}